\documentclass[a4paper,12pt]{article}

\title{PISNO OCENJEVANJE ZNANJA -- ZAPOREDJA}
\author{}
\date{}

\begin{document}

\maketitle

\begin{center}
	\large\textbf{Čas reševanja: 60 minut}
\end{center}

\vspace{0.5cm}

\noindent
\begin{minipage}[t]{0.54\textwidth}
	\textbf{Kriterij ocenjevanja}

	\medskip

	\begin{tabular}{|c|c|}
		\hline
		Ocena & Odstotek uspešnosti \\
		\hline
		1 & 0--49\% \\
		\hline
		2 & 50--64\% \\
		\hline
		3 & 65--79\% \\
		\hline
		4 & 80--89\% \\
		\hline
		5 & 90--100\% \\
		\hline
	\end{tabular}

	\medskip

	\begin{tabular}{ll}
		\textbf{Dosežene točke:} & \rule{1.5cm}{0.4pt} \\
		\textbf{Možno število točk:} & 30 \\
		\textbf{Ocena:} & \rule{1.5cm}{0.4pt}
	\end{tabular}
\end{minipage}
\hfill
\begin{minipage}[t]{0.38\textwidth}
	\textbf{Podatki dijaka}

	\vspace{0.4cm}

	\textbf{Ime:}\\[0.2cm]
	\rule{\linewidth}{0.4pt}\\[0.7cm]
	\textbf{Priimek:}\\[0.2cm]
	\rule{\linewidth}{0.4pt}\\[0.7cm]
	\textbf{Razred:}\\[0.2cm]
	\rule{\linewidth}{0.4pt}\\[0.7cm]
	\textbf{Datum:}\\[0.2cm]
	\rule{\linewidth}{0.4pt}
\end{minipage}

\newpage

\section*{Naloge}

\noindent\textit{Pri vsaki nalogi zapiši postopek reševanja.}

\begin{enumerate}
	\item \textbf{(6 točk)} Aritmetično zaporedje ima $a_3=11$ in $a_8=31$.
	\begin{enumerate}
		\item Izračunaj diferenco $d$ in prvi člen $a_1$.
		\item Zapiši splošni člen zaporedja in izračunaj $a_{25}$.
		\item Izračunaj vsoto prvih 25 členov.
	\end{enumerate}
	\newpage

	\item \textbf{(6 točk)} Geometrijsko zaporedje ima $a_2=6$, $a_5=162$ in količnik $q>0$.
	\begin{enumerate}
		\item Izračunaj količnik $q$ in prvi člen $a_1$.
		\item Zapiši splošni člen zaporedja.
		\item Izračunaj vsoto prvih šestih členov.
	\end{enumerate}
	\newpage

	\item \textbf{(6 točk)} Zaporedje je podano s splošnim členom
	\[
		a_n=\frac{2n+1}{n+2}, \qquad n=1,2,\ldots
	\]
	\begin{enumerate}
		\item Zapiši prve štiri člene zaporedja.
		\item Ugotovi, ali je zaporedje naraščajoče, in utemelji odgovor. Določi tudi njegovo zgornjo mejo.
		\item Izračunaj limito zaporedja.
	\end{enumerate}
	\newpage

	\item \textbf{(6 točk)} Zaporedje je določeno rekurzivno:
	\[
		a_1=4, \qquad a_{n+1}=a_n+2n+1.
	\]
	\begin{enumerate}
		\item Izračunaj prvih pet členov.
		\item Zapiši domnevo za splošni člen $a_n$.
		\item Domnevo dokaži z matematično indukcijo.
	\end{enumerate}
	\newpage

	\item \textbf{(6 točk)} V gledališču je v prvi vrsti 18 sedežev, v vsaki naslednji vrsti pa so trije sedeži več kot v prejšnji. Gledališče ima 20 vrst.
	\begin{enumerate}
		\item Koliko sedežev je v zadnji vrsti?
		\item Koliko sedežev je v gledališču skupaj?
	\end{enumerate}
\end{enumerate}

\end{document}
